\chapter{Einleitung}

\section{Motivation und Problemstellung}

Scheduling-Verfahren gehören zu den zentralen Konzepten moderner Betriebssysteme. Sie bestimmen, welcher Prozess zu welchem Zeitpunkt CPU-Zeit erhält und beeinflussen unmittelbar Reaktionsfähigkeit, Fairness und Auslastung eines Systems. Obwohl diese Verfahren in Lehrbüchern und Vorlesungen theoretisch sauber beschrieben werden, fällt es Studierenden häufig schwer, ihren zeitlichen Ablauf in konkreten Situationen zu verstehen \cite{silberschatz2018operating,tanenbaum2014modern}.

Gerade präemptive Verfahren sind für Lernende anspruchsvoll, weil sie viele Zustandswechsel, Unterbrechungen und Neuplanungen erzeugen. Die Entscheidungslogik bleibt bei reinen Textbeschreibungen und statischen Beispielen oft abstrakt. Das Problem zeigt sich besonders beim Vergleich verschiedener Verfahren unter identischen Eingaben: Die fachlichen Unterschiede sind zwar theoretisch bekannt, aber in ihrer unmittelbaren Wirkung auf Ablauf, Zeitverlauf und Fairness nicht immer intuitiv nachvollziehbar.

Vor diesem Hintergrund wurde eine interaktive Webanwendung entwickelt, die Scheduling-Verfahren als sichtbaren, steuerbaren Ablauf modelliert. Ziel ist nicht die Beantwortung einer rein wissenschaftlichen Forschungsfrage, sondern die Erstellung eines nutzbaren, didaktischen Werkzeugs, das technische Zusammenhänge transparent macht und Studierenden einen praktischen Zugang zu den Grundlagen des CPU-Scheduling eröffnet.

\section{Zielsetzung der Arbeit}

Die vorliegende Arbeit verfolgt die Konzeption und Umsetzung einer interaktiven Webanwendung zur Veranschaulichung präemptiver Scheduling-Verfahren. Das Ziel besteht darin, eine Anwendung zu entwickeln, die:

\begin{itemize}
    \item ausgewählte Scheduling-Verfahren fachlich korrekt simuliert,
    \item ihre zeitlichen Abläufe verständlich visualisiert,
    \item verschiedene Verfahren anhand identischer Szenarien vergleichbar macht,
    \item und als didaktisches Werkzeug im Lehrkontext nutzbar ist.
\end{itemize}

Der Fokus liegt damit auf der praktischen Umsetzung eines verständlichen und erweitbaren Produkts. Die Arbeit verbindet fachliche Grundlagen aus der Betriebssystemtheorie mit Prinzipien aus Software Engineering und Human-Computer Interaction. Dabei steht die nachvollziehbare Architektur, die klare Interaktion und die unmittelbare Produktwirkung im Vordergrund.

\section{Thematische Ausrichtung der Arbeit}

Die Arbeit ist bewusst nicht als klassische empirische Untersuchung konzipiert, sondern als engineering-orientierte Produktdokumentation. Der Schwerpunkt liegt auf der Gestaltung einer Anwendung, die im Studienkontext einen Mehrwert bietet und zugleich technische, didaktische und gestalterische Entscheidungen nachvollziehbar dokumentiert.

Diese Ausrichtung ist für den Studiengang geeignet, weil sie die Verbindung zwischen theoretischem Wissen und praktischer Umsetzung sichtbar macht. Die fachlichen Grundlagen aus der Betriebssystemlehre werden mit einer strukturierten Architektur, benutzerfreundlicher Oberfläche und konkreten Produktbeispielen verknüpft. Dadurch entsteht eine Bachelorarbeit, die nicht nur theoretisch fundiert ist, sondern auch die Entwicklung eines nutzbaren Softwaresystems als zentrale Leistung hervortellt.

\section{Vorgehensweise und Aufbau der Arbeit}

Die Arbeit folgt einem engineering-orientierten Vorgehen:

\begin{enumerate}
    \item Analyse der fachlichen Grundlagen zu CPU-Scheduling und relevanten Messgrößen,
    \item Ableitung von Anforderungen an eine didaktische Visualisierungsanwendung,
    \item Konzeption der Systemarchitektur und der Komponentenstruktur,
    \item technische Umsetzung der Webanwendung,
    \item produktorientierte Darstellung der Ergebnisse mit Screenshots und systematischer Erklärung.
\end{enumerate}

Der Aufbau der Bachelorarbeit ist entsprechend auf die Entwicklung und Präsentation des Produkts ausgerichtet. Kapitel 2 behandelt die fachlichen Grundlagen. Kapitel 3 und Kapitel 4 legen die Anforderungen an die Architektur und die HCI-gestützte Gestaltung fest. Kapitel 5 beschreibt die konkrete Konzeption und den Entwurf der Anwendung. Kapitel 6 dokumentiert die technische Implementierung. Kapitel 7 stellt das Ergebnis mit der Produktpräsentation und der Visualisierung der wichtigsten Funktionen dar. Kapitel 8 schließt mit Reflexion, Grenzen und Ausblick.

\section{Leitgedanke der Arbeit}

Der zentrale Gedanke dieser Arbeit lautet:

\begin{quote}
    Praktische Verständlichkeit entsteht nicht nur durch fachliche Korrektheit, sondern auch durch eine gut strukturierte Architektur, klare Interaktion und eine verständliche Darstellung des zeitlichen Ablaufs.
\end{quote}

Diese Perspektive bildet den Kern der Arbeit. Sie macht deutlich, dass die Simulation von Scheduling-Verfahren nicht nur eine fachliche Aufgabe ist, sondern auch ein Design- und Architekturproblem mit unmittelbaren Auswirkungen auf Nutzbarkeit und Lernwirkung.
